\subsection{Real-Model Judgments of Occurrence--Parameter Relationships}
\label{subsec:relationship-audit}
Retrieving a suitable schema leaves open which source occurrence expresses each parameter. We therefore examine whether actual LLM services give complete, structurally grounded and mutually consistent judgments on proposed occurrence--declaration relationships. This is a separate annotation/QC question from both the 331-query scalar-grounding study and the common-18 formulation comparison. It tests neither the autonomous correctness of the retrieval system nor accuracy against human labels.

The frozen October 27, 2026 records contain 217 distinct relationship items from 16 NL4Opt training problems \cite{ramamonjison2023nl4opt}, drawn as a category-covering, high-risk convenience packet from a previously inspected 100-problem review frame. They are not 217 independent problems, repeated independent trials, or a random corpus sample. Multiple relationships share problem text, numerical occurrences and declarations. Each service originally received 33 batches of up to eight targets, with complete source text, supplied variables, objective and constraints; peer judgments were withheld. The four nominal decisions are \texttt{CORRECT}, \texttt{INCORRECT}, \texttt{AMBIGUOUS} and \texttt{CANNOT\_DETERMINE}. These are model assertions, not verified truth labels. Each object must also supply an explanation, literal source quotations, valid source identifiers, omission notes, a shared-mention flag and uncertainty.

The requested models are \path{gemini-3.1-pro-preview} through Google Vertex AI, \path{deepseek-ai/DeepSeek-V4-Pro-0813} through io.net io.intelligence (recorded FP8), and \path{command-a-reasoning-08-2025} through Cohere's direct v2 API. DeepSeek here names the served model, not a direct DeepSeek API. All use temperature 0 and a 12,000-output-token cap; Gemini uses LOW thinking, Cohere a 256-token thinking budget, and io.net its default reasoning behavior. Sampling/analysis seed 20261027 is not an inference seed. The identifiers authenticate the requested or echoed service names, not underlying weight checksums or every server default.

\begin{table}[htbp]\centering\small
\caption{Structural validity on the same 217 relationship targets per service. Final counts combine original-valid objects, deterministic offline repairs and historically generated valid replacements under a revised format. They are not single-pass or equal-treatment accuracy estimates. No new calls were made for this revision.}\label{tab:relationship-validity}
\setlength{\tabcolsep}{3.5pt}
\begin{tabular}{@{}lrrrr@{}}\toprule
Service/model family & Original valid & Offline repairs & Valid replacements & Final valid\\\midrule
Vertex AI/Gemini & 215 & 2 & 0 & 217/217\\
io.net/DeepSeek & 180 & 33 & 4 & 217/217\\
Cohere/Command A & 105 & 24 & 68 & 197/217\\\midrule
Total & 500 & 59 & 72 & 631/651\\\bottomrule
\end{tabular}\end{table}

Strict output-contract checks require exact source quotations and membership in supplied variable/formulation identifiers, in addition to JSON structure. Recovery v2.1 preserves all 500 originally valid objects, extracts only unique complete envelopes, maps only explicitly documented source aliases and losslessly serializes structured omission notes. Of 92 model-item targets still invalid after offline repair, historical smaller-batch requests with explicit identifier/quotation instructions yield 72 valid replacements. Cohere additionally uses a native JSON Schema with identifier enums; an initial rejected schema and a corrected transport are retained in the lineage. These changes affect comparability. The remaining 20 invalid Cohere objects all contain nonliteral evidence quotations; invalidity does not prove a mathematically wrong decision.

Among jointly valid pairs, nominal agreement is Gemini--DeepSeek $213/217$ ($98.16\%$), Gemini--Cohere $123/197$ ($62.44\%$), and DeepSeek--Cohere $120/197$ ($60.91\%$). The complete three-model denominator is 197, comprising 119 determinate unanimous cases, 77 disagreements and one unresolved case. The remaining 20 of 217 are invalid/incomplete. The unresolved case has three \texttt{CORRECT} labels but no verified lexical endpoint, so raw label unanimity is $120/197$, while determinate unanimity is $119/197$. Neither is semantic accuracy.

Fleiss' nominal $\kappa$ on the 197 complete cases is $0.370807$ (observed pair agreement $0.739425$, chance term $0.585858$), using all four decision categories and excluding invalid objects rather than coding them as an opinion. The complete-case marginal counts are 420 \texttt{CORRECT}, 168 \texttt{INCORRECT}, two \texttt{CANNOT\_DETERMINE} and one \texttt{AMBIGUOUS}, totaling 591 judgments. Original complete cases number only 89, with $\kappa=0.501643$; those same 89 cases retain identical labels after recovery. The change in aggregate $\kappa$ reflects expanded complete-case membership, not measured loss or gain in correctness. We report no population confidence interval, significance claim or model ranking for these clustered, selected items.

Verifiable defects include a singular variable alias rejected against a plural source identifier, prose/structured-note output that violates the schema, and quotations that paraphrase rather than copy the source. Valid decisions also differ on source-occurrence binding: Supplement, Section~\ref{S-app:relationship-evidence} gives identified witnesses and the per-problem denominator ledger. These examples explain protocol failures or conflicting model assertions, not independently adjudicated semantic error frequencies. No human-verified reference judgments exist for the 217 targets, and no consensus is approved as training gold. The study complements the retrieval/readiness gap: a retrieved template and a plausible structured judgment can both look usable while correct source binding remains unverified. Agreement cannot validate the retrieval pipeline's autonomous performance.
